\chapter{FUNCTIONAL ANALYSIS AND MODELING}

\section{Use Case Modeling}

A use case model shows what a system does for its users, without going into how it is built \cite{uml_systems_analysis}. SmartExam has 12 use cases, one for each module M1 to M12. The same 12 use cases, with the same numbers and names, are used for the user stories and use case tables in this chapter, the activity and sequence diagrams in Chapter~5 and the unit test cases in Chapter~7.

\DiagramFig{width=\linewidth,height=0.8\textheight,keepaspectratio}{use_case_diagram.png}%
{Use Case Diagram of SmartExam}{fig:usecase_diagram}%
{Figure~\ref{fig:usecase_diagram} shows the five actors and the 12 use cases. Every user signs in (UC-01); the student then uses UC-02 to UC-06 and the staff use UC-07 to UC-12.}

\subsection{User Stories}

A user story states what a user wants and why \cite{user_stories_applied}. The stories below are grouped by actor, and each one names the use case it belongs to.

\subsubsection{Student}
\begin{itemize}[nosep]
    \item As a student, I want to log in on my lab PC and have my account bound to it, so that nobody can sit my exam from another machine (UC-01).
    \item As a student, I want to run a readiness check and a practice exam before the real exam, so that problems with my PC are found early (UC-02).
    \item As a student, I want to see my assigned exam, its countdown and my eligibility, and start it when it opens, so that I know whether and when I can sit it (UC-03).
    \item As a student, I want to see my results, violation history and notifications on the web portal after the exam, so that I can follow my progress (UC-03).
    \item As a student, I want to sit the exam in a locked fullscreen window, so that every candidate works under the same rules (UC-04).
    \item As a student, I want my PC to send its status and screenshots during the exam, so that the invigilator can see that I am online and working (UC-05).
    \item As a student, I want my answers saved automatically and a receipt when I submit, so that my work is not lost if the network or the PC fails (UC-06).
\end{itemize}

\subsubsection{Teacher}
\begin{itemize}[nosep]
    \item As a teacher, I want to sign in to the web panel with my email and password, so that only I can work on my exams (UC-01).
    \item As a teacher, I want to create an exam with MCQ, coding and theory questions and set its lockdown rules, so that the exam runs under the rules I chose (UC-08).
    \item As a teacher, I want to send an announcement about an exam, also by email if I choose, so that every student knows the schedule (UC-08).
    \item As a teacher, I want to import attendance, set eligibility, override single students and assign seats and invigilators, so that only eligible students sit the exam at their own PCs (UC-09).
    \item As a teacher, I want to watch every student live and warn, add time, unlock the PC or force a submission, so that I can deal with problems during the exam (UC-10).
    \item As a teacher, I want coding answers run against test cases and theory answers given suggested marks, so that marking is faster while I keep the final mark (UC-11).
    \item As a teacher, I want to see how similar the submissions are, so that I can check likely copying myself (UC-11).
    \item As a teacher, I want to download an exam's reports and evidence, so that I have a record of the exam (UC-12).
\end{itemize}

\subsubsection{Invigilator}
\begin{itemize}[nosep]
    \item As an invigilator, I want to sign in to the web panel, so that I can open the exams I am assigned to (UC-01).
    \item As an invigilator, I want to see the live status of the exams I am assigned to, so that I can watch every student in the lab (UC-10).
    \item As an invigilator, I want to warn a student, add time, unlock the PC, record an incident or force a submission, so that I can act at once when something goes wrong (UC-10).
\end{itemize}

\subsubsection{Admin and SuperAdmin}
\begin{itemize}[nosep]
    \item As an admin, I want to sign in to the web panel, and to be refused after a wrong password, so that nobody else can use my account (UC-01).
    \item As an admin, I want to create, import, deactivate and sign out users and set their roles, so that each person has only the access they need (UC-07).
    \item As an admin, I want to keep departments and courses and reset a student's device binding when a PC is replaced, so that the records stay correct (UC-07).
    \item As a SuperAdmin, I want to be the only one who manages Admin and SuperAdmin accounts, so that administrator rights cannot be given out freely (UC-07).
    \item As an admin, I want to register the labs and their PCs, so that each student's seat can be checked when the exam starts (UC-09).
    \item As an admin, I want to read the audit log, check that no record was changed and see system analytics, so that I can tell who did what (UC-12).
\end{itemize}

\subsection{Individual Actor Use Cases}

% Each lead-in paragraph is kept on the same page as the use case table after it.
\makeatletter
\newsavebox{\uctbox}
\newcommand{\uctlead}[1]{\par{\interlinepenalty\@M\clubpenalty\@M\widowpenalty\@M #1\par}\nopagebreak}
\makeatother

Tables~\ref{tab:uc01} to~\ref{tab:uc12} describe the 12 use cases in the same fields \cite{use_case_definition}. The Steps row gives the normal flow, and the Exception Flow row gives what can go wrong and what the system then does.

\begin{lrbox}{\uctbox}{\small\renewcommand{\arraystretch}{1.1}\setlength{\tabcolsep}{6pt}%
\begin{tabular}{>{\raggedright\arraybackslash}p{3.6cm}>{\raggedright\arraybackslash}p{\dimexpr\textwidth-3.6cm-4\tabcolsep\relax}}
\hline\hline
\textbf{Field} & \textbf{Description} \\ \hline
UCT ID & UC-01 \\
Case Name & Log In and Bind Device \\
Use-Case Overview & A user signs in with a password on the web panel. For an exam, the student signs in on the desktop client, and the account is also bound to one lab PC. \\
Actors & Student, Teacher, Invigilator, Admin, SuperAdmin \\
Pre-Conditions & The user has an active account; for an exam, the desktop client runs on a lab PC. \\
Triggers & The user presses Sign In. \\
Steps & 1.~The user enters an email (a student on the desktop client may use the roll number) and the password.\newline
2.~The server checks the password.\newline
3.~On the web panel, the server opens a session and the user's dashboard opens.\newline
4.~On the desktop client, an ID of the PC hardware is also sent. At the first login the server binds the account to this PC; at every later login the PC must match.\newline
5.~The server opens the student's session, and the exam dashboard opens. \\
Post-Conditions & The user is signed in; on the desktop client, a student only on the bound PC. \\
Exception Flow & Wrong password or inactive account: the login is refused and recorded. Too many login attempts in one minute: further attempts are refused for a while. A student on another PC: the login is refused until an admin resets the binding (UC-07). \\
\hline\hline
\end{tabular}}\end{lrbox}
\uctlead{Table~\ref{tab:uc01} is about signing in. Every user signs in with a password, and a student's account is also bound to the lab PC at the first login.}
\begin{table}[H]
\centering
\caption{Use Case Table for Log In and Bind Device}\label{tab:uc01}
\vspace{4pt}
\usebox{\uctbox}
\end{table}

\begin{lrbox}{\uctbox}{\small\renewcommand{\arraystretch}{1.1}\setlength{\tabcolsep}{6pt}%
\begin{tabular}{>{\raggedright\arraybackslash}p{3.6cm}>{\raggedright\arraybackslash}p{\dimexpr\textwidth-3.6cm-4\tabcolsep\relax}}
\hline\hline
\textbf{Field} & \textbf{Description} \\ \hline
UCT ID & UC-02 \\
Case Name & Run Readiness Check \\
Use-Case Overview & The client checks that the PC can run the exam; the student can also try a practice exam that does not count. \\
Actors & Student \\
Pre-Conditions & The student is signed in (UC-01). \\
Triggers & The student presses Run check or Practice on the dashboard. \\
Steps & 1.~The client checks the network, the clock, screen capture, the webcam and microphone if the exam needs them, a single monitor with no remote session or virtual machine, and free disk space.\newline
2.~The client shows the result of each check.\newline
3.~The server checks the result again and stores it. \\
Post-Conditions & The result is stored and shown to the staff on the live monitor. \\
Exception Flow & A check fails: its reason is shown, and the student fixes it and runs the check again. If the exam requires readiness, a pass older than 30~minutes must be repeated. \\
\hline\hline
\end{tabular}}\end{lrbox}
\uctlead{Table~\ref{tab:uc02} covers the readiness check and the practice exam. The check runs on the PC before Start is allowed, and the server keeps the result.}
\begin{table}[H]
\centering
\caption{Use Case Table for Run Readiness Check}\label{tab:uc02}
\vspace{4pt}
\usebox{\uctbox}
\end{table}

\begin{lrbox}{\uctbox}{\small\renewcommand{\arraystretch}{1.1}\setlength{\tabcolsep}{6pt}%
\begin{tabular}{>{\raggedright\arraybackslash}p{3.6cm}>{\raggedright\arraybackslash}p{\dimexpr\textwidth-3.6cm-4\tabcolsep\relax}}
\hline\hline
\textbf{Field} & \textbf{Description} \\ \hline
UCT ID & UC-03 \\
Case Name & View Exam Dashboard and Start Exam \\
Use-Case Overview & The dashboard shows the student's assigned exam, and the student checks in and starts it. \\
Actors & Student \\
Pre-Conditions & The student is signed in (UC-01). \\
Triggers & The student opens the dashboard after logging in. \\
Steps & 1.~The dashboard shows the assigned exam, its instructions, the countdown and the eligibility, readiness and seat status.\newline
2.~The student presses Check In.\newline
3.~Start Exam is enabled at the start time, after check-in and, if the exam requires it, a readiness pass (UC-02).\newline
4.~The student presses Start Exam; the server checks eligibility, the seat, the exam time and readiness, and opens the exam. \\
Post-Conditions & The exam session is in progress and the exam window opens (UC-04). After the exam, the student sees the results, violation history, performance and notifications on the web portal. \\
Exception Flow & Ineligible student: Start Exam stays disabled and the reason is shown. Wrong seat, exam not open yet, exam closed or readiness missing: the server refuses the start and shows the reason. \\
\hline\hline
\end{tabular}}\end{lrbox}
\uctlead{Table~\ref{tab:uc03} is the student's dashboard. Start Exam opens the exam only when the student is eligible, at the right seat and inside the exam time.}
\begin{table}[H]
\centering
\caption{Use Case Table for View Exam Dashboard and Start Exam}\label{tab:uc03}
\vspace{4pt}
\usebox{\uctbox}
\end{table}

\begin{lrbox}{\uctbox}{\small\renewcommand{\arraystretch}{1.1}\setlength{\tabcolsep}{6pt}%
\begin{tabular}{>{\raggedright\arraybackslash}p{3.6cm}>{\raggedright\arraybackslash}p{\dimexpr\textwidth-3.6cm-4\tabcolsep\relax}}
\hline\hline
\textbf{Field} & \textbf{Description} \\ \hline
UCT ID & UC-04 \\
Case Name & Sit the Locked Exam \\
Use-Case Overview & The student answers the questions in a locked fullscreen window. It always includes UC-05. \\
Actors & Student \\
Pre-Conditions & The exam session is in progress (UC-03). \\
Triggers & The exam starts, or an interrupted exam is resumed. \\
Steps & 1.~The exam window opens fullscreen and stays on top.\newline
2.~Alt+Tab, the Windows key and Task Manager are blocked.\newline
3.~The client checks the running apps every 2~s against the exam's allowed apps, and watches the clipboard, USB drives, extra monitors, remote sessions and virtual machines.\newline
4.~Each breach is reported to the server (UC-05).\newline
5.~The student answers the questions and submits (UC-06). \\
Post-Conditions & The exam is submitted and the PC is released. \\
Exception Flow & App that is not allowed: in Report mode it is reported, in Enforce mode it is also closed. A warning, extra time or unlock from the invigilator (UC-10) appears in the window; after an unlock the window locks itself again. \\
\hline\hline
\end{tabular}}\end{lrbox}
\uctlead{Table~\ref{tab:uc04} describes the locked exam window. The PC stays locked until the exam is submitted, and the exam's lockdown mode decides what happens to an app that is not allowed.}
\begin{table}[H]
\centering
\caption{Use Case Table for Sit the Locked Exam}\label{tab:uc04}
\vspace{4pt}
\usebox{\uctbox}
\end{table}

\begin{lrbox}{\uctbox}{\small\renewcommand{\arraystretch}{1.1}\setlength{\tabcolsep}{6pt}%
\begin{tabular}{>{\raggedright\arraybackslash}p{3.6cm}>{\raggedright\arraybackslash}p{\dimexpr\textwidth-3.6cm-4\tabcolsep\relax}}
\hline\hline
\textbf{Field} & \textbf{Description} \\ \hline
UCT ID & UC-05 \\
Case Name & Send Monitoring Data \\
Use-Case Overview & While the exam is open, the client sends the PC's status, events, screenshots and webcam snapshots to the server. \\
Actors & Student (through the client; included in UC-04) \\
Pre-Conditions & The locked exam is open (UC-04). \\
Triggers & The exam window opens. \\
Steps & 1.~Every 10~s the client sends a heartbeat with focus, running apps, idle time and camera status.\newline
2.~Events are sent in small batches every 5~s; the server marks each one as a warning, suspicious or a violation.\newline
3.~Screenshots are taken at random times (every 30 to 90~s by default) and at once after a violation.\newline
4.~If the exam needs the webcam, the client checks each frame on the PC and sends a snapshot when a problem is found.\newline
5.~The server stores the data and shows it on the live monitor. \\
Post-Conditions & Events, screenshots and snapshots are stored as evidence for the exam. \\
Exception Flow & Network lost: the data waits in a local queue and is sent later. No heartbeat for 30~s: the server marks the student as Disconnected. \\
\hline\hline
\end{tabular}}\end{lrbox}
\uctlead{Table~\ref{tab:uc05} describes the monitoring data that the client sends during the exam. The invigilator sees this data on the live monitor (UC-10) and in the reports (UC-12).}
\begin{table}[H]
\centering
\caption{Use Case Table for Send Monitoring Data}\label{tab:uc05}
\vspace{4pt}
\usebox{\uctbox}
\end{table}

\begin{lrbox}{\uctbox}{\small\renewcommand{\arraystretch}{1.1}\setlength{\tabcolsep}{6pt}%
\begin{tabular}{>{\raggedright\arraybackslash}p{3.6cm}>{\raggedright\arraybackslash}p{\dimexpr\textwidth-3.6cm-4\tabcolsep\relax}}
\hline\hline
\textbf{Field} & \textbf{Description} \\ \hline
UCT ID & UC-06 \\
Case Name & Autosave and Submit Answers \\
Use-Case Overview & Answers are saved during the exam and submitted with a receipt. \\
Actors & Student \\
Pre-Conditions & The exam session is in progress (UC-04). \\
Triggers & The student edits an answer or presses Submit Exam, or the time runs out. \\
Steps & 1.~Each edit is saved to an encrypted backup on the PC.\newline
2.~The answer is saved to the server 10~s after the edit, and at once when the student changes question.\newline
3.~The student presses Submit Exam and confirms.\newline
4.~The server closes the session and returns a receipt of the stored answers, and the answers wait for marking (UC-11). \\
Post-Conditions & The exam is submitted or timed out, and its receipt is stored. \\
Exception Flow & Network lost: failed saves are kept and sent again in order. PC crash: the answers are restored from the server and the backup. Time runs out: the client submits by itself. \\
\hline\hline
\end{tabular}}\end{lrbox}
\uctlead{Table~\ref{tab:uc06} covers saving and submitting answers. Each answer is kept on the PC and on the server, so a network or PC failure does not lose work.}
\begin{table}[H]
\centering
\caption{Use Case Table for Autosave and Submit Answers}\label{tab:uc06}
\vspace{4pt}
\usebox{\uctbox}
\end{table}

\begin{lrbox}{\uctbox}{\small\renewcommand{\arraystretch}{1.1}\setlength{\tabcolsep}{6pt}%
\begin{tabular}{>{\raggedright\arraybackslash}p{3.6cm}>{\raggedright\arraybackslash}p{\dimexpr\textwidth-3.6cm-4\tabcolsep\relax}}
\hline\hline
\textbf{Field} & \textbf{Description} \\ \hline
UCT ID & UC-07 \\
Case Name & Manage Users and Roles \\
Use-Case Overview & An administrator keeps the user accounts, their roles, the departments and courses and the device bindings. \\
Actors & Admin, SuperAdmin \\
Pre-Conditions & The admin is signed in on the web panel. \\
Triggers & The admin opens the Users, Departments \& Courses or Device Bindings page. \\
Steps & 1.~The admin creates or edits a user and sets the role.\newline
2.~The admin imports users from a CSV file; each bad row is skipped and listed with its reason.\newline
3.~The admin deactivates, activates or deletes a user, or signs a user out at once.\newline
4.~The admin adds departments and courses, and resets a student's device binding when a PC is replaced.\newline
5.~Every change is recorded in the audit log. \\
Post-Conditions & The account is stored, and its role decides what the user may do. \\
Exception Flow & Email already in use or unknown role: the change is refused. An Admin tries to create or change an Admin or SuperAdmin account: refused and recorded, because only a SuperAdmin may do this. \\
\hline\hline
\end{tabular}}\end{lrbox}
\uctlead{Table~\ref{tab:uc07} covers the accounts and the academic records kept by an administrator. Every change is recorded in the audit log, and only a SuperAdmin manages Admin and SuperAdmin accounts.}
\begin{table}[H]
\centering
\caption{Use Case Table for Manage Users and Roles}\label{tab:uc07}
\vspace{4pt}
\usebox{\uctbox}
\end{table}

\begin{lrbox}{\uctbox}{\small\renewcommand{\arraystretch}{1.1}\setlength{\tabcolsep}{6pt}%
\begin{tabular}{>{\raggedright\arraybackslash}p{3.6cm}>{\raggedright\arraybackslash}p{\dimexpr\textwidth-3.6cm-4\tabcolsep\relax}}
\hline\hline
\textbf{Field} & \textbf{Description} \\ \hline
UCT ID & UC-08 \\
Case Name & Create and Configure Exam \\
Use-Case Overview & The teacher defines an exam, its questions and its lockdown rules, and announces it. \\
Actors & Teacher \\
Pre-Conditions & The teacher is signed in, and the course section exists. \\
Triggers & The teacher opens Create Exam. \\
Steps & 1.~The teacher enters the title, start time, duration and instructions, and selects the students.\newline
2.~The teacher adds MCQ, coding and theory questions with their marks, test cases and model answers.\newline
3.~The teacher sets the allowed apps, the lockdown mode (Off, Report or Enforce), readiness, webcam and microphone options and the screenshot interval.\newline
4.~The server checks and stores the exam.\newline
5.~The teacher sends an announcement to the students, also by email if chosen. \\
Post-Conditions & The exam is scheduled, and its students see it on their dashboard (UC-03). \\
Exception Flow & A missing or invalid setting: the server refuses to save and shows the error. The teacher can edit the exam later. \\
\hline\hline
\end{tabular}}\end{lrbox}
\uctlead{Table~\ref{tab:uc08} describes how a teacher sets up an exam. The settings chosen here are the rules that UC-02 to UC-06 enforce on the lab PC.}
\begin{table}[H]
\centering
\caption{Use Case Table for Create and Configure Exam}\label{tab:uc08}
\vspace{4pt}
\usebox{\uctbox}
\end{table}

\begin{lrbox}{\uctbox}{\small\renewcommand{\arraystretch}{1.1}\setlength{\tabcolsep}{6pt}%
\begin{tabular}{>{\raggedright\arraybackslash}p{3.6cm}>{\raggedright\arraybackslash}p{\dimexpr\textwidth-3.6cm-4\tabcolsep\relax}}
\hline\hline
\textbf{Field} & \textbf{Description} \\ \hline
UCT ID & UC-09 \\
Case Name & Manage Eligibility and Seats \\
Use-Case Overview & The teacher sets who is eligible from attendance and assigns seats and invigilators; the admin registers labs and PCs. \\
Actors & Teacher, Admin \\
Pre-Conditions & The exam exists and has students (UC-08). \\
Triggers & The teacher opens the Eligibility page, or the admin opens the Labs \& Workstations page. \\
Steps & 1.~The admin registers labs and their PCs with name and IP address.\newline
2.~The teacher uploads the section's attendance CSV file.\newline
3.~The teacher recalculates eligibility; students below the minimum attendance (75\% by default) become ineligible.\newline
4.~The teacher overrides single students where needed, with a note.\newline
5.~The teacher assigns a PC to each student and invigilators to the exam.\newline
6.~Every import, recalculation and override is recorded in the audit log. \\
Post-Conditions & Each student has an eligibility and, where set, a seat; an ineligible student cannot start (UC-03). \\
Exception Flow & Missing or bad file: refused, and bad rows are listed by line. Recalculation after the exam has ended: refused. Overridden students are kept as set when eligibility is recalculated. \\
\hline\hline
\end{tabular}}\end{lrbox}
\uctlead{Table~\ref{tab:uc09} decides who may sit the exam and where. The teacher works from attendance, and the admin registers the labs and PCs that the seats use.}
\begin{table}[H]
\centering
\caption{Use Case Table for Manage Eligibility and Seats}\label{tab:uc09}
\vspace{4pt}
\usebox{\uctbox}
\end{table}

\begin{lrbox}{\uctbox}{\small\renewcommand{\arraystretch}{1.1}\setlength{\tabcolsep}{6pt}%
\begin{tabular}{>{\raggedright\arraybackslash}p{3.6cm}>{\raggedright\arraybackslash}p{\dimexpr\textwidth-3.6cm-4\tabcolsep\relax}}
\hline\hline
\textbf{Field} & \textbf{Description} \\ \hline
UCT ID & UC-10 \\
Case Name & Supervise Exam Live \\
Use-Case Overview & Staff watch every student of a running exam and act on problems. \\
Actors & Invigilator, Teacher \\
Pre-Conditions & The exam is running, and the invigilator is assigned to it or the teacher owns it. \\
Triggers & The staff member opens the Live Monitor and picks an exam. \\
Steps & 1.~A tile for each student updates live and shows the status: Stable, Warning, Suspicious or Critical.\newline
2.~The staff member opens a student to see the timeline, screenshots and webcam snapshots.\newline
3.~The staff member sends a warning, adds time, unlocks the PC for a short time, records an incident or force-submits the exam.\newline
4.~Each action is recorded in the audit log. \\
Post-Conditions & The student sees the warning, extra time or unlock on the exam window (UC-04). \\
Exception Flow & No heartbeat for 30~s: the tile shows Disconnected until the PC reconnects. A staff member who is not assigned to the exam: the action is refused. \\
\hline\hline
\end{tabular}}\end{lrbox}
\uctlead{Table~\ref{tab:uc10} is the live monitor used by the invigilator and the teacher. Each action reaches the student's PC with its next heartbeat, within about 10 seconds.}
\begin{table}[H]
\centering
\caption{Use Case Table for Supervise Exam Live}\label{tab:uc10}
\vspace{4pt}
\usebox{\uctbox}
\end{table}

\begin{lrbox}{\uctbox}{\small\renewcommand{\arraystretch}{1.1}\setlength{\tabcolsep}{6pt}%
\begin{tabular}{>{\raggedright\arraybackslash}p{3.6cm}>{\raggedright\arraybackslash}p{\dimexpr\textwidth-3.6cm-4\tabcolsep\relax}}
\hline\hline
\textbf{Field} & \textbf{Description} \\ \hline
UCT ID & UC-11 \\
Case Name & Evaluate Answers \\
Use-Case Overview & The system marks the answers and checks them for copying; the teacher reviews the suggested marks. \\
Actors & Teacher \\
Pre-Conditions & The exam sessions are submitted (UC-06). \\
Triggers & The teacher opens the results of an exam. \\
Steps & 1.~MCQ answers are marked against the answer key.\newline
2.~C and C++ answers are compiled and run against the weighted test cases.\newline
3.~Theory answers get a suggested mark from TF-IDF similarity with the model answer \cite{salton1988}.\newline
4.~Answers are compared in pairs by k-gram winnowing \cite{winnowing}, and pairs at or above the exam's threshold (70\% by default) are flagged.\newline
5.~The teacher accepts the suggested marks or overrides a mark with a note. \\
Post-Conditions & The suggested marks and the teacher's marks are stored separately; flagged pairs are only listed, with no automatic penalty. \\
Exception Flow & Compile error or time-out: the test case fails and the message is shown. No code runner available: the answer is left for the teacher to mark. \\
\hline\hline
\end{tabular}}\end{lrbox}
\uctlead{Table~\ref{tab:uc11} covers marking after the exam. The system only suggests marks, and the teacher's decision is the final mark.}
\begin{table}[H]
\centering
\caption{Use Case Table for Evaluate Answers}\label{tab:uc11}
\vspace{4pt}
\usebox{\uctbox}
\end{table}

\begin{lrbox}{\uctbox}{\small\renewcommand{\arraystretch}{1.1}\setlength{\tabcolsep}{6pt}%
\begin{tabular}{>{\raggedright\arraybackslash}p{3.6cm}>{\raggedright\arraybackslash}p{\dimexpr\textwidth-3.6cm-4\tabcolsep\relax}}
\hline\hline
\textbf{Field} & \textbf{Description} \\ \hline
UCT ID & UC-12 \\
Case Name & View Reports and Audit Log \\
Use-Case Overview & Staff get an exam's reports and evidence, and an admin checks the audit log. \\
Actors & Teacher, Admin \\
Pre-Conditions & The staff member is signed in; for reports, the teacher owns the exam. \\
Triggers & The teacher opens the Reports page, or the admin opens the Audit Logs page. \\
Steps & 1.~The teacher picks a report: the exam summary, attendance, violations or scores as PDF, an Excel workbook, or an evidence ZIP file.\newline
2.~The server builds the file and adds it to the exam's report list.\newline
3.~The teacher downloads the file.\newline
4.~The admin filters the audit log and runs the check that shows whether any record was changed.\newline
5.~The admin views the system analytics. \\
Post-Conditions & The report is stored and can be downloaded again; the audit check result is shown. \\
Exception Flow & Unknown report type or no exam chosen: refused. Teacher does not own the exam: refused. Audit check fails: the first changed record is shown. \\
\hline\hline
\end{tabular}}\end{lrbox}
\uctlead{Table~\ref{tab:uc12} covers the records of an exam. The teacher downloads reports and evidence, and the admin checks the audit log and the system analytics.}
\begin{table}[H]
\centering
\caption{Use Case Table for View Reports and Audit Log}\label{tab:uc12}
\vspace{4pt}
\usebox{\uctbox}
\end{table}
